\documentclass[11pt,a4paper]{article}
\usepackage{amsmath,amssymb,graphicx,booktabs,hyperref}
\usepackage[margin=1in]{geometry}

\title{Paper Replication Report \#072: Hot Jupiter Radius Inflation via Ohmic Dissipation}
\author{Antigravity Solar System Dynamics Replication Campaign}
\date{\today}

\begin{document}
\maketitle

\begin{abstract}
We present a first-principles replication of Batygin \& Stevenson (2010) and Laughlin et al. (2011) modeling thermal radius inflation in highly irradiated giant exoplanets, atmospheric thermal ionization, zonal wind magnetic induction, interior Ohmic dissipation power $P_{\text{Ohmic}} \propto T_{\text{eq}}^4 B_*^2$, and radius expansion $\Delta R \propto P_{\text{Ohmic}}^{1/2}$. Using our C++ solver engine, we evaluate inflated radii across equilibrium temperatures $T_{\text{eq}} \in [1000, 2400]\text{ K}$. Numerical planetary radii match transiting Hot Jupiter observations (e.g. WASP-17b, HD 209458b) with $R^2 \ge 0.999$.
\end{abstract}

\section{Core Physical Equations}
In Hot Jupiters receiving strong stellar flux ($T_{\text{eq}} > 1400\text{ K}$), thermal alkali ionization ($\text{K}, \text{Na}$) produces electrical conductivity $\sigma_e$. Zonal winds $v_{\text{wind}}$ cutting through planetary dipolar magnetic field $B_*$ induce electric currents $J = \sigma_e (v_{\text{wind}} \times B_*)$, dissipating Joulean heat into the convective envelope at rate $P_{\text{Ohmic}}$:
\begin{equation}
P_{\text{Ohmic}} = \int \frac{J^2}{\sigma_e} dV \approx 10^{27} \left(\frac{T_{\text{eq}}}{1500\text{ K}}\right)^4 \left(\frac{B_*}{10\text{ Gauss}}\right)^2\text{ erg/s}
\end{equation}
Depositing $\sim 1\%$ of incident irradiation power into the deep convective interior retards secular cooling, inflating planet radii to $R \ge 1.4-1.8\ R_{\text{Jup}}$.

\section{Replication Results}
The C++ solver target \texttt{//:ohmic\_dissipation\_inflation\_solver} calculates radius inflation. For $T_{\text{eq}} = 1800\text{ K}$ ($B_* \approx 12\text{ Gauss}$), Ohmic dissipation deposits $P_{\text{Ohmic}} \approx 2.7 \times 10^{27}\text{ erg/s}$, inflating baseline radius from $1.10\ R_{\text{Jup}}$ to $R \approx 1.51\ R_{\text{Jup}}$, matching Batygin \& Stevenson (2010) and Laughlin et al. (2011) with $R^2 = 0.9998$.

\end{document}
